\chapter{Opis funkcjonalny systemu}

Poniżej przedstawione zostały najważniejsze funkcje zapewniane przez aplikację \verb|Wydatkomat|. Należy mieć na uwadze, że aplikacja desktopowa została stworzona głównie jako panel administratora, w związku z czym nie zapewnia wszystkich funkcji. Z drugiej strony aplikacja mobilna i webowa nie zapewniają możliwości administrowania użytkownikami.

\section{Uwierzytelnianie i autoryzacja użytkowników}

Aby korzystać z aplikacji należy się zarejestrować. Jest to pierwsza funkcjonalność, która jest obligatoryjna, aby w pełni korzystać z systemu. Użytkownik może zarejestrować się formularzem zaimplementowanym w klientach lub poprzez protokół OAuth2 (Google, Facebook). Obie metody wiążą się z założeniem konta w systemie.

Po rejestracji należy się zalogować, sposoby logowania są analogiczne do sposobów rejestracji.

Użytkownik może połączyć swoje lokalnie utworzone konto z wieloma mediami społecznościowymi za pomocą protokołu OAuth2, co pozwoli na logowanie się obiema metodami.

Z drugiej strony użytkownik, który założył konto logując się przez OAuth2, ma możliwość utworzenia hasła, aby móc logować się formularzem z poziomu aplikacji.

Po poprawnym zalogowaniu przypisane zostają do użytkownika dwa tokeny: \verb|JWT| oraz \verb|refresh token|. Token dostępowy JWT jest ważny przez okres 1 godziny, po tym czasie sesja zostaje unieważniona. Aby ponownie móc korzystać z systemu należy się zalogować lub wykorzystać refresh token, który jest ważny przez tydzień.

Dodatkowo system umożliwia włączenie logowania dwuetapowego przy użyciu aplikacji mobilnej \verb|Google Authenticator|. Aby to zrobić należy uruchomić opcję z poziomu profilu. Po włączeniu tej funkcji i poprawnym połączeniu aplikacji z Wydatkomatem, przy kolejnych próbach logowania należy przejść również drugi etap logowania -- zweryfikowanie się kodem z aplikacji Google Authenticator.

\section{Zarządzanie wydatkami}

Główną funkcją aplikacji jest tworzenie nowych wydatków i edycja już istniejących. Użytkownik może utworzyć wydatek, dodając do niego swoich znajomych. Wydatek może mieć tytuł, opis, kwotę i wielu uczestników. Tylko osoba tworząca wydatek może go edytować. Przy edycji wydatków bez problemu można zmieniać tytuł i opis. Można również zmieniać kwotę i uczestników, lecz w tym przypadku system waliduje aktualnie spłacone kwoty i może zablokować taką edycję. Jeśli edycja zostanie dozwolona, następuje rekalkulacja należnych kwot.

W aktualnej wersji systemu dostępne jest jedynie dzielenie wydatków po równo, co oznacza, że w przypadku utworzenia wydatku na 500 złotych i dodaniu 4 znajomych, kwota zostanie podzielona na 125 zł / osobę. Właściciel wydatku nie może być jego uczestnikiem.

Wydatki oraz ich uczestników można również usuwać.

\section{Znajomi}

Uczestnikami wydatków mogą być jedynie znajomi osoby, która go utworzyła. Każdy użytkownik może dodać znajomego po wyszukaniu go z listy aktywnych użytkowników systemu.

Po wysłaniu zaproszenia do znajomych osoba zapraszana otrzymuje stosowne powiadomienie, które może zaakceptować lub odrzucić.

Użytkownicy mogą sprawdzić oczekujące zaproszenia, zarówno wysłane jak i odebrane. Mogą również zobaczyć listę swoich znajomych.

Istnieje możliwość zakończenia znajomości z innym użytkownikiem. Znajomości są dwukierunkowe, co oznacza, że po akceptacji zaproszenia do znajomych obie strony widzą siebie na liście znajomych i mogą tworzyć między sobą wydatki.

\section{Opłaty}

Uczestnicy wydatku widzą sumaryczną kwotę wydatku, swoją część do spłaty oraz część, którą już spłacili.

Wydatek można opłacać stopniowo, czyli do jednego wydatku pojedynczy użytkownik może stworzyć wiele spłat.

Istnieje możliwość usunięcia stworzonej przez siebie płatności.

\section{Profil}

Zalogowany użytkownik może podejrzeć i edytować swój profil.

Z poziomu profilu można zarządzać swoim kontem, uwzględniając w to:
\begin{itemize}
	\item ustawienie hasła lokalnego (jeśli takie jeszcze nie istnieje, dotyczy użytkowników OAuth2),
	\item połączenie kont społecznościowych,
	\item zmianę hasła,
	\item włączenie/wyłączenie logowania dwuetapowego,
	\item aktualizację swoich danych osobowych,
	\item usunięcie konta.
\end{itemize}

\section{Powiadomienia}

Każdy użytkownik może sterować kanałami powiadomień. W aplikacji na obecnym etapie dostępne są dwa kanały -- email i websocket (powiadomienia realtime w aplikacji).

Powiadomienia dotyczą np.:
\begin{itemize}
	\item zaproszenia do znajomych,
	\item nowego wydatku,
	\item nieudanej próby logowania.
\end{itemize}

Powiadomienia związane z bezpieczeństwem są na stałe przypisane do kanału \verb|email|.

Preferencje dotyczące powiadomień są edytowalne. Powiadomienia można oznaczyć jako przeczytane w przypadku powiadomień websocket.

\section{Zarządzanie użytkownikami}

Jako administrator istnieje możliwość zarządzania użytkownikami.

Administrator widzi pełną listę aktywnych użytkowników, może dezaktywować użytkownika, zresetować lub zmienić jego hasło oraz utworzyć nowego użytkownika.

\section{Statystyki}

System umożliwia podgląd okresowych statystyk zarówno dotyczących wydatków jak i spłat. Umożliwia to lepsze zrozumienie i kontrolowanie budżetu.

Statystyki uwzględniają:
\begin{itemize}
	\item całkowitą liczbę wydatków/spłat,
	\item całkowitą kwotę wydatków/spłat,
	\item średnią kwotę wydatków/spłat.
\end{itemize}